% W4i29_P_updates.tex
% DFD 17-Agent Research Programme — Iteration 29
% Agents 2 (Lab/MEMS), 3 (Clocks), 5 (Lensing/GW), 6 (P1-P4), 7 (P5-P7), 8 (P8-P11), 9 (mochi_class), 13 (Particle)
% Date: 2026-03-21

\documentclass[11pt,a4paper]{article}
\usepackage[utf8]{inputenc}
\usepackage{amsmath,amssymb,amsfonts}
\usepackage{hyperref}
\usepackage{booktabs}
\usepackage{longtable}
\usepackage{geometry}
\geometry{margin=2.5cm}

\title{DFD i29 Patent \& Experimental Updates\\Agents 2, 3, 5, 6, 7, 8, 9, 13}
\author{DFD 17-Agent Research Programme}
\date{2026-03-21}

\begin{document}
\maketitle
\tableofcontents
\newpage

%=============================================================================
\section{Agent 6: Patents P1--P4 --- Dead Patents Final Check}
%=============================================================================

\subsection{Mandate}
Final check on P1--P4 (earliest DFD patents). Any resurrections in 2025--2026?

\subsection{New Literature (i29)}

\begin{enumerate}
\item \textbf{``The Speed of Gravity and the Fate of Dark Energy''} --- arXiv: \href{https://arxiv.org/abs/2511.19762}{2511.19762} (2025).\\
Review of post-GW170817 constraints on scalar-tensor gravity. $c_T = c$ confirmed to $\mathcal{O}(10^{-15})$. No mechanism to resurrect theories with $c_T \neq c$. P1--P4's original formulations remain dead. \textbf{Grade: B.}

\item \textbf{``Recent Developments in Degenerate Higher Order Scalar Tensor Theories''} --- arXiv: \href{https://arxiv.org/abs/2407.18234}{2407.18234} (2024, updated 2025).\\
DHOST classification review. Confirms that only Class Ia (where DFD sits) survives GW170817 + decay constraints. No loopholes found. \textbf{Grade: B.}

\item \textbf{Patent landscape search (2025--2026):} No new filings or grants related to P1--P4 concepts. No resurrections identified. \textbf{Grade: N/A.}
\end{enumerate}

\subsection{Verdict}
\textbf{P1--P4: CONFIRMED DEAD.} No resurrection pathway exists. The GW170817 constraint is fundamental and cannot be evaded without abandoning the core assumptions of P1--P4. All viable DFD physics now lives in P5--P15 and the core axiom system.

\textbf{RECOMMENDATION (i29):} Remove P1--P4 from active monitoring after i30. Final iteration of surveillance. Resources better allocated to P5--P15 and new predictions.

\textbf{Grade: RED (dead). Stable since i5.}

%=============================================================================
\section{Agent 7: Patents P5--P7 --- ET Governance}
%=============================================================================

\subsection{Mandate}
Track Einstein Telescope ERIC governance decision and site selection. Update P5--P7 timelines.

\subsection{New Literature (i29)}

\begin{enumerate}
\item \textbf{ET bid book deadline:} December 2026 (confirmed). Board of Governmental Representatives (BGR) site decision: \textbf{second half of 2027}. Construction start: 2028. First light: 2035. \textbf{Grade: A (milestone update).}

\item \textbf{Flanders increases ET reserve to \EUR{500}M} --- (September 2025). \href{https://www.einsteintelescope-emr.eu/en/2025/09/22/flanders-increases-einstein-telescope-reserve-to-500-million-euros/}{Link}.\\
Major financial commitment from Belgian region. Strengthens Meuse--Rhine Euroregion (EMR) site candidacy. \textbf{Grade: B.}

\item \textbf{Sardinia Region allocates additional \EUR{20}M for ET} --- (April 2025). \href{https://www.einstein-telescope.it/en/2025/04/07/the-sardinia-region-allocates-another-e20-million-for-the-einstein-telescope/}{Link}.\\
Sardinia (Sos Enattos mine site) continues investment. Site competition between EMR and Sardinia remains active. \textbf{Grade: B.}

\item \textbf{CTAO ERIC precedent} --- (2025). \href{https://www.eurekalert.org/news-releases/1069629}{EurekAlert link}.\\
Cherenkov Telescope Array Observatory became an ERIC in 2025, establishing governance precedent for ET. \textbf{Grade: C.}
\end{enumerate}

\subsection{ET Timeline for DFD}
\begin{center}
\begin{tabular}{lll}
\toprule
\textbf{Date} & \textbf{Milestone} & \textbf{DFD Relevance} \\
\midrule
Dec 2026 & Bid book deadline & --- \\
H2 2027 & BGR site decision & --- \\
2028 & Construction start & --- \\
2035 & ET first light & GW speed, lensing, polarization tests \\
2035+ & ET science operations & $\Sigma(y)$ via GW lensing statistics \\
\bottomrule
\end{tabular}
\end{center}

\textbf{Assessment:} ET remains the ultimate GW test for DFD but is 9+ years away. Near-term GW constraints come from LVK O5 (starting 2026--2027). P5--P7 are long-horizon patents.

\textbf{Grade: GREEN. On track but distant.}

%=============================================================================
\section{Agent 8: Patents P8--P11 --- SQMS 2.0}
%=============================================================================

\subsection{Mandate}
Track SQMS 2.0 renewal status and proposal mechanics. Update Dark SRF timelines for P8--P11.

\subsection{New Literature (i29)}

\begin{enumerate}
\item \textbf{SQMS 2.0 renewal: \$125M over 5 years} --- Announced November 4, 2025. \href{https://news.fnal.gov/2025/11/fermilabs-sqms-center-funded-with-125-million-to-shape-the-future-of-quantum-information-science/}{Fermilab news}.\\
\textbf{CRITICAL UPDATE.} SQMS has been renewed. \$25M first year, 36 partners. Key goals: 100-qudit processor, 10 ms transmon coherence, quantum data centre infrastructure. Dark SRF searches continue under the renewed mandate. \textbf{Grade: A.}

\item \textbf{DOE NQI renewal: \$625M across 5 centres} --- November 2025. \href{https://www.rigzone.com/news/usa_energy_dept_announces_625mm_to_renew_5_nqisr_centers-11-nov-2025-182294-article/}{Link}.\\
All five National Quantum Information Science Research Centers renewed. Total: \$625M. \textbf{Grade: A.}

\item \textbf{Dark SRF widely-tunable cavity} --- Cervantes et al.\ (2025). \href{https://indico.global/event/14966/contributions/134111/attachments/63287/122128/RaphaelCervantesCPAD2025.pdf}{CPAD 2025 talk}.\\
New widely-tunable SRF cavity for dark photon searches covering broader mass range. Design sensitivity $g_{a\gamma\gamma} \sim 2 \times 10^{-11}$ GeV$^{-1}$. \textbf{Grade: B.}

\item \textbf{UMN Dark Photon Hunt under SQMS renewal} --- \href{https://cse.umn.edu/physics/news/does-quantum-center-sqms-renewal-fuels-umns-hunt-dark-photons}{UMN Physics}.\\
University of Minnesota group specifically funded for dark photon searches under SQMS 2.0. SRF cavities at design sensitivity expected 2026--2027. \textbf{Grade: B.}

\item \textbf{Fermilab quantum sensor study (March 2026)} --- \href{https://news.fnal.gov/2026/03/fermilab-led-study-advances-development-of-sophisticated-quantum-sensors-to-track-high-energy-particles-and-detect-dark-matter/}{Fermilab news}.\\
Advances in quantum sensors for dark matter detection. SRF-based approaches highlighted. \textbf{Grade: B.}
\end{enumerate}

\subsection{SQMS 2.0 --- Impact on DFD}

\textbf{The renewal changes the landscape for P8--P11.} The original ``Phase 0'' proposal (\$700K/30 months) was designed as a \emph{pilot} within SQMS 1.0. With SQMS 2.0 now funded:

\begin{enumerate}
\item \textbf{Dark SRF infrastructure is secured for 5 years.} No need for a separate funding proposal for cavity access.
\item \textbf{The UMN dark photon group is specifically funded.} Contact: UMN Physics (SQMS collaboration).
\item \textbf{Design sensitivity reached 2026--2027.} The widely-tunable cavity covers a broader mass range than the original Phase 0 concept.
\item \textbf{Proposal pathway:} Instead of a standalone Phase 0, DFD-specific search parameters can be proposed as a \emph{parasitic search mode} within the existing SQMS dark photon programme. This requires:
\begin{itemize}
\item Identifying the specific frequency/mass window where DFD's scalar-field coupling would produce a signal
\item Writing a 2--3 page white paper for the SQMS Internal Review Committee
\item Contact: SQMS Center Director (Fermilab)
\end{itemize}
\end{enumerate}

\textbf{CORRECTION from i28:} The \$700K/30-month Phase 0 proposal is now \textbf{superseded} by the SQMS 2.0 renewal. The path forward is integration into the existing programme, not a standalone proposal.

\textbf{Grade: GREEN. Upgraded from YELLOW.} SQMS 2.0 provides the infrastructure. DFD needs to write a white paper for parasitic search mode.

%=============================================================================
\section{Agent 2: Lab/MEMS --- MAGIS-100 and Atom Interferometry}
%=============================================================================

\subsection{Mandate}
Track MAGIS-100 commissioning status and other atom interferometry experiments.

\subsection{New Literature (i29)}

\begin{enumerate}
\item \textbf{MAGIS-100 laser lab construction complete} --- January 2026. \href{https://news.fnal.gov/2026/01/fermilab-completes-laser-lab-construction-for-worlds-largest-vertical-atom-interferometer/}{Fermilab news}.\\
\textbf{Major milestone.} Laser infrastructure for the 100m vertical atom interferometer complete. Next: atom source delivery (late 2026), vacuum tube installation (2027), commissioning (2028). \textbf{Grade: A.}

\item \textbf{Temples --- MAGIS-100 at CETUP 2025} --- \href{https://lss.fnal.gov/archive/2025/slides/fermilab-slides-25-0118-etd.pdf}{Fermilab slides}.\\
Environmental characterization ongoing. Kovachy (Northwestern) leading vibration/noise studies. Laser beam trajectory fluctuations from ground vibrations being characterized. \textbf{Grade: B.}

\item \textbf{3rd Terrestrial Very-Long-Baseline Atom Interferometry Workshop} --- August 2025, CERN. \href{https://indico.cern.ch/event/1522117/}{Indico link}.\\
International coordination meeting including MAGIS-100, ZAIGA, AION, MIGA. Harmonized sensitivity forecasts for gravitational wave detection and dark matter searches. \textbf{Grade: B.}

\item \textbf{ZAIGA status:} Phase 1 (300m vertical tunnel) under construction near Wuhan. Full ZAIGA-GW (3 $\times$ 1 km triangular) target: 2035. Current status: early R\&D. Not competitive with MAGIS-100 timeline. \textbf{Grade: C.}
\end{enumerate}

\subsection{MAGIS-100 Timeline}
\begin{center}
\begin{tabular}{lll}
\toprule
\textbf{Date} & \textbf{Milestone} & \textbf{Status} \\
\midrule
Jan 2026 & Laser lab complete & DONE \\
Late 2026 & Atom sources delivered (Stanford) & On track \\
2027 & Vacuum tube + atom source installation & Planned \\
End 2027 & Installation complete, early data & Planned \\
2028 & Full commissioning & Planned \\
2028+ & Science operations & --- \\
\bottomrule
\end{tabular}
\end{center}

\textbf{DFD Relevance:} MAGIS-100 is sensitive to scalar-field gradients in the $0.1$--$10$ Hz band. DFD's $\psi$ field produces a local gravitational gradient $a = (c^2/2)\nabla\psi$. The sensitivity to $\nabla\psi$ depends on the baseline length and atom temperature. At 100m baseline, MAGIS-100 can probe $|\nabla\psi| \gtrsim 10^{-15}$ m$^{-1}$, corresponding to DFD effects at $\sim 10^3 a_0$.

\textbf{Grade: GREEN. On track for 2028 commissioning.}

%=============================================================================
\section{Agent 3: Clocks --- Th-229 Nuclear Clock}
%=============================================================================

\subsection{Mandate}
Track Th-229 nuclear clock progress. Update P29 timeline.

\subsection{New Literature (i29)}

\begin{enumerate}
\item \textbf{``Frequency reproducibility of solid-state thorium-229 nuclear clocks''} --- Nature (2025). \href{https://www.nature.com/articles/s41586-025-09999-5}{DOI link}.\\
\textbf{Major result.} Demonstrates reproducibility of the $^{229}$Th nuclear clock transition in CaF$_2$ crystals across doping concentrations, temperatures, and time. Establishes solid-state nuclear clocks as viable frequency standards. \textbf{Grade: A.}

\item \textbf{``Frequency ratio of the $^{229m}$Th nuclear isomeric transition and the $^{87}$Sr atomic clock''} --- Nature (2024). \href{https://www.nature.com/articles/s41586-024-07839-6}{DOI link}.\\
First absolute frequency measurement of the Th-229 isomer transition referenced to the Sr lattice clock. Establishes the connection between nuclear and atomic clock standards. \textbf{Grade: A.}

\item \textbf{``Fine-structure constant sensitivity of the Th-229 nuclear clock transition''} --- Nature Communications (2025). \href{https://www.nature.com/articles/s41467-025-64191-7}{DOI link}.\\
Measures $K = 5900 \pm 2300$ for the Th-229 transition's sensitivity to $\alpha$ variations. Confirms enhanced sensitivity over atomic clocks by factor $\sim 10^4$. \textbf{Critical for DFD.} \textbf{Grade: A.}

\item \textbf{``Toward a Thorium-229 Isomer-Based Nuclear Clock: Recent Progress''} --- Nuclear Physics News 35(2), 2025. \href{https://www.tandfonline.com/doi/full/10.1080/10619127.2025.2483617}{DOI link}.\\
Review of state of the art. Five active groups: NIST, PTB, TU Wien, UCLA, JILA. Competition accelerating. \textbf{Grade: B.}

\item \textbf{UCLA December 2025 breakthrough:} Hudson group achieved laser-driven internal conversion electron M\"{o}ssbauer spectroscopy of $^{229}$ThO$_2$. Breaks reliance on VUV-transparent host crystals. Expands material options. \textbf{Grade: B.}
\end{enumerate}

\subsection{P29 Status Update}

The P29 prediction stated that transportable Th-229 nuclear clocks would achieve $\delta\nu/\nu \sim 10^{-18}$ sensitivity to fundamental constant variations by 2027--2028. The i29 literature shows:

\begin{itemize}
\item Solid-state reproducibility demonstrated (Nature 2025): ON TRACK
\item $\alpha$-sensitivity $K = 5900 \pm 2300$ confirmed: ON TRACK (and $K$ value supports DFD's prediction of enhanced sensitivity)
\item Five competing groups: TIMELINE ACCELERATING
\item UCLA ThO$_2$ approach removes VUV crystal bottleneck: EXPANDS REACH
\end{itemize}

\textbf{DFD connection:} The scalar field $\psi$ couples to the fine-structure constant as $\alpha_\mathrm{eff} = \alpha_0 e^{d_e \psi}$ where $d_e$ is the scalar-electromagnetic coupling. The Th-229 nuclear clock's $K \approx 5900$ amplification means it probes $\psi$-induced $\alpha$ variations at the level $\delta\alpha/\alpha \sim K \cdot \delta\nu/\nu \sim 5900 \times 10^{-18} \sim 6 \times 10^{-15}$. This is competitive with Equivalence Principle tests.

\textbf{Grade: GREEN. P29 on track for 2027--2028.}

%=============================================================================
\section{Agent 5: Lensing/GW --- O4 Lensing Candidates}
%=============================================================================

\subsection{Mandate}
Track LVK O4 lensing search results. Update GW lensing statistics for DFD.

\subsection{New Literature (i29)}

\begin{enumerate}
\item \textbf{GWTC-4.0 released} --- August 2025. \href{https://www.ligo.caltech.edu/news/ligo20250826}{LIGO news}.\\
128 new GW signal candidates from O4. Total GWTC catalogue now $>$300 events. Companion papers include lensing search. \textbf{Grade: A.}

\item \textbf{LVK O4 completed} --- November 18, 2025. \href{https://www.ligo.caltech.edu/news/ligo20251118}{LIGO news}.\\
O4 ran from May 2023 to November 2025. O4a + O4b + O4c data now in analysis pipeline. \textbf{Grade: A (milestone).}

\item \textbf{``Lensing, not luck! Detection prospects of strongly lensed gravitational waves''} --- arXiv: \href{https://arxiv.org/abs/2510.23238}{2510.23238} (2025).\\
First $3\sigma$ detection of strong GW lensing expected in O5, not O4. O4 provides upper limits on lensing rate. \textbf{Grade: B.}

\item \textbf{O5 planning:} Six-month observing run beginning late summer/early fall 2026. Detectors participating as available. LIGO sensitivity upgrade underway. \textbf{Grade: B.}

\item \textbf{LVK March 2026 catalogue expansion:} 128 new candidates more than doubling total detections. \href{https://phys.org/news/2026-03-lvk-gravitational-candidates.html}{Phys.org}.\\
\textbf{Grade: B.}
\end{enumerate}

\subsection{DFD GW Lensing Assessment}

DFD predicts that gravitational lensing of GWs is modified by the disformal coupling $\Sigma(y)$. Specifically, the lensing magnification of GWs differs from electromagnetic lensing by:
\begin{equation}
\mu_\mathrm{GW} / \mu_\mathrm{EM} = \Sigma(y)^{-2}
\end{equation}

For $y = 0.15$: $\mu_\mathrm{GW} / \mu_\mathrm{EM} \approx 1.16$. This 16\% enhancement is too small to detect with individual O4 events but could be constrained statistically with O5 (which is expected to find the first lensed GW events).

\textbf{Grade: GREEN-YELLOW. O4 provides no lensed events; O5 (2026--2027) is the next opportunity.}

%=============================================================================
\section{Agent 9: mochi\_class --- SymBoltz.jl Deep Dive}
%=============================================================================

\subsection{Mandate}
Assess SymBoltz.jl as an alternative to modifying CLASS for DFD Boltzmann computation. Compare feasibility, learning curve, and timeline.

\subsection{New Literature (i29)}

\begin{enumerate}
\item \textbf{``SymBoltz.jl: A symbolic-numeric, approximation-free, and differentiable linear Einstein--Boltzmann solver''} --- Hersle (2026). A\&A, March 2026. arXiv: \href{https://arxiv.org/abs/2509.24740}{2509.24740}.\\
\textbf{KEY PAPER.} Julia-based Boltzmann solver. Symbolic equation specification, no approximation switching, automatic differentiation compatible. Published in A\&A March 2026. \textbf{Grade: A.}

\item \textbf{SymBoltz.jl GitHub repository} --- \href{https://github.com/hersle/SymBoltz.jl}{github.com/hersle/SymBoltz.jl}.\\
Open source. MIT license. Active development. Documentation and tutorials available. \textbf{Grade: A (tool).}

\item \textbf{``DISCO-DJ I: a differentiable Einstein-Boltzmann solver for cosmology''} --- arXiv: \href{https://arxiv.org/abs/2311.03291}{2311.03291} (2023, updated 2025).\\
Python-based alternative. Differentiable but slower than SymBoltz.jl. Uses JAX backend. \textbf{Grade: B.}

\item \textbf{CosmoCentral.jl Boltzmann Solver} --- \href{https://www.marcobonici.com/CosmoCentral.jl/dev/BoltzmannSolver/}{Documentation}.\\
Another Julia-based cosmology toolkit with Boltzmann solver capabilities. Less mature than SymBoltz.jl. \textbf{Grade: C.}
\end{enumerate}

\subsection{SymBoltz.jl vs mochi\_class: Feasibility Comparison}

\begin{center}
\begin{longtable}{p{4cm}p{5cm}p{5cm}}
\toprule
\textbf{Criterion} & \textbf{mochi\_class (CLASS mod)} & \textbf{SymBoltz.jl} \\
\midrule
\endfirsthead
\textbf{Language} & C (CLASS core) + Python wrapper & Julia \\
\textbf{Equation input} & Hard-coded; modify C source & Symbolic; equations specified as math \\
\textbf{Approximations} & Tight-coupling, RSA, UFA switches & None; implicit solver handles stiffness \\
\textbf{Autodiff} & No & Yes (forward/reverse mode) \\
\textbf{DFD modifications} & 5-step plan, 8 months (i26 estimate) & Potentially 2--3 months \\
\textbf{Learning curve} & CLASS internals are complex; 2000+ line modifications & Julia + ModelingToolkit.jl; moderate \\
\textbf{Validation} & Extensive (CLASS community) & Validated against CLASS to $0.1\%$ \\
\textbf{Community} & Large & Small (single developer + early adopters) \\
\textbf{MCMC integration} & CosmoMC, MontePython, Cobaya & Custom (but autodiff enables gradient-based sampling) \\
\textbf{Risk} & High (C debugging, approximation interactions) & Medium (less community support) \\
\bottomrule
\end{longtable}
\end{center}

\subsection{SymBoltz.jl Advantages for DFD}

\begin{enumerate}
\item \textbf{Symbolic equation input:} DFD's modified Friedmann equations, perturbation equations with $\mu(k,z)$ and $\Sigma(y)$, and the scalar field evolution $\ddot{\psi} + 3H\dot{\psi} = -\partial V/\partial\psi$ can be entered directly as symbolic equations. No need to navigate CLASS's C code architecture.

\item \textbf{No approximation switching:} The implicit solver handles the stiff equations at all times. This eliminates the risk of DFD modifications interacting badly with CLASS's tight-coupling and radiation streaming approximations --- which was the \emph{primary technical risk} of the mochi\_class approach.

\item \textbf{Automatic differentiation:} Enables gradient-based MCMC (e.g., HMC, NUTS) for parameter estimation. This could be 10--100$\times$ faster than Metropolis-Hastings for DFD's multi-parameter space ($A_s$, $n_s$, $\Omega_m$, $H_0$, $y_0$, $W''(0)$, ...).

\item \textbf{Development speed:} Estimated 2--3 months for a DFD implementation vs 8 months for CLASS modification.
\end{enumerate}

\subsection{SymBoltz.jl Risks}

\begin{enumerate}
\item \textbf{Single developer:} The package is maintained by one person (H.~Hersle). Bus factor = 1.
\item \textbf{Community validation:} Less extensively tested than CLASS. Potential for undiscovered bugs.
\item \textbf{Performance:} Julia's JIT compilation is fast, but SymBoltz.jl's implicit solver may be slower per evaluation than CLASS's optimized explicit solver. Mitigated by autodiff enabling fewer total evaluations.
\item \textbf{MCMC integration:} No off-the-shelf connection to standard cosmological samplers. Would require custom sampling pipeline.
\end{enumerate}

\subsection{Recommendation}

\textbf{SWITCH PRIMARY APPROACH TO SymBoltz.jl.} The advantages (symbolic input, no approximation switching, autodiff, faster development) outweigh the risks (single developer, smaller community). The mochi\_class 5-step/8-month plan should be retained as a \textbf{fallback} for validation.

\textbf{Proposed timeline:}
\begin{enumerate}
\item Month 1: Install SymBoltz.jl, reproduce $\Lambda$CDM $C_\ell$ spectra, validate against CLASS.
\item Month 2: Implement DFD background ($n(\psi)$, modified Friedmann) and perturbation ($\mu$, $\Sigma$) equations.
\item Month 3: Generate DFD $C_\ell^{TT}$, $C_\ell^{EE}$, $C_\ell^{\phi\phi}$. Compare to Planck. Assess CMB third-peak deficit.
\item Month 4 (optional): Set up gradient-based MCMC for DFD parameter estimation.
\end{enumerate}

\textbf{Grade: GREEN. Major strategic shift.}

%=============================================================================
\section{Agent 13: Particle Physics --- JUNO First Results}
%=============================================================================

\subsection{Mandate}
Track JUNO commissioning and first oscillation results.

\subsection{New Literature (i29)}

\begin{enumerate}
\item \textbf{``First measurement of reactor neutrino oscillations at JUNO''} --- JUNO Collaboration (2025). arXiv: \href{https://arxiv.org/abs/2511.14593}{2511.14593}.\\
\textbf{MAJOR RESULT.} Physics data taking began August 26, 2025. From 59.1 days of data: $\sin^2\theta_{12} = 0.3092 \pm 0.0087$, $\Delta m^2_{21} = (7.50 \pm 0.12) \times 10^{-5}$ eV$^2$. Factor 1.6$\times$ better precision than all previous experiments combined. \textbf{Grade: A.}

\item \textbf{``Neutrinoless Double Beta Decay in Light of JUNO First Data''} --- arXiv: \href{https://arxiv.org/abs/2511.15391}{2511.15391} (2025).\\
Implications of JUNO's precise $\theta_{12}$ measurement for $0\nu\beta\beta$ searches. Tightens the allowed parameter space. \textbf{Grade: B.}

\item \textbf{JUNO ``solar neutrino tension'' confirmation:} First data confirms the mild $1.5\sigma$ discrepancy between solar and reactor determinations of $\Delta m^2_{21}$. JUNO will test this with both channels. \textbf{Grade: B.}
\end{enumerate}

\subsection{DFD Relevance}

JUNO's primary DFD relevance is through the neutrino mass hierarchy determination (expected with 3--4 years of data, i.e., 2028--2029). DFD predicts specific relationships between the scalar field $\psi$ and effective neutrino masses through the conformal coupling. The mass hierarchy determines the sign of $\Delta m^2_{31}$, which enters the neutrino contribution to the matter power spectrum.

Secondary relevance: the ``solar neutrino tension'' ($1.5\sigma$ discrepancy in $\Delta m^2_{21}$ between solar and reactor neutrinos) could, if confirmed, indicate new physics in neutrino propagation. DFD's scalar field modifies the effective potential experienced by neutrinos in matter (MSW effect modified by $\psi$). However, the current $1.5\sigma$ level is not significant.

\textbf{Grade: GREEN. JUNO operational and producing world-leading data.}

%=============================================================================
\section{Patent \& Experimental Status Summary}
%=============================================================================

\begin{center}
\begin{longtable}{llll}
\toprule
\textbf{Agent} & \textbf{Topic} & \textbf{Status} & \textbf{Change from i28} \\
\midrule
\endfirsthead
6 & P1--P4 (dead patents) & RED (dead) & No change \\
7 & P5--P7 (ET/GW) & GREEN & Site decision H2 2027 (was ``2026'') \\
8 & P8--P11 (SQMS/Dark SRF) & GREEN & Upgraded; SQMS 2.0 funded \$125M \\
2 & MAGIS-100 & GREEN & Laser lab complete \\
3 & Th-229 nuclear clock & GREEN & Solid-state reproducibility in Nature \\
5 & GW lensing (O4) & GREEN-YELLOW & O4 complete, no lensed events \\
9 & mochi\_class/SymBoltz.jl & GREEN & \textbf{Strategic shift to SymBoltz.jl} \\
13 & JUNO & GREEN & First results published \\
\bottomrule
\end{longtable}
\end{center}

\textbf{TOP FINDING (Agent 8):} SQMS 2.0 has been renewed at \$125M/5yr. The standalone Phase 0 proposal is superseded. DFD should pursue parasitic search mode via SQMS white paper.

\textbf{TOP FINDING (Agent 9):} SymBoltz.jl (published A\&A March 2026) offers a dramatically faster path to DFD Boltzmann computation than modifying CLASS. Recommended as primary approach with mochi\_class as fallback.

\end{document}
